\documentclass[10pt]{article}

\usepackage{hyperref}
\usepackage{xcolor}
\usepackage{calc}
\usepackage{graphicx}
\usepackage{tikz}
\usepackage{fontspec}
\usepackage{fontawesome5}
\usepackage{titlesec}
\usepackage{enumitem}
\usepackage{fancybox}
\usepackage{xeCJK}

\hypersetup{hidelinks}
\urlstyle{same}

%%%%%%%%%%%%%%%%%%%%
% 设置
%%%%%%%%%%%%%%%%%%%%

\setlength{\parindent}{0pt}
\setlength{\parskip}{0pt}
\pagenumbering{gobble}
\setlist[itemize]{
    leftmargin=1.35em,
    itemsep=0.12em,
    topsep=0.18em,
    parsep=0pt,
    partopsep=0pt
}
\setlist[enumerate]{leftmargin=*}
\renewcommand{\arraystretch}{1.12}
\linespread{1.08}

\titleformat{\section}
  {\Large\bfseries\raggedright}
  {}{0em}
  {}
  [{\color{secondarycolor}\titlerule}]
\titlespacing*{\section}{0pt}{0.55em}{0.28em}

\usepackage[
    a4paper,
    left=1.1cm,
    right=1.1cm,
    top=0.75cm,
    bottom=0.85cm,
    nohead
]{geometry}

\setCJKmainfont[
    Path=fonts/,
    Extension=.otf,
    BoldFont=*-Bold,
]{NotoSerifSC}

\definecolor{primarycolor}{RGB}{200, 60, 60}
\definecolor{secondarycolor}{RGB}{200, 68, 28}

\newlength{\iconwidth}
\setlength{\iconwidth}{1.5em}

\newcommand{\projectheading}[3]{%
    \noindent\makebox[\textwidth][s]{%
        \makebox[0pt][l]{{\large \textbf{#1}}}%
        \hfill
        \makebox[0pt][c]{{\normalsize \textbf{#2}}}%
        \hfill
        \makebox[0pt][r]{#3}%
    }%
}

%%%%%%%%%%%%%%%%%%%%
% 文章内容
%%%%%%%%%%%%%%%%%%%%

\newcommand{\contact}{
    \footnotesize
    \textcolor{white}{
        \faEnvelope \quad 邮箱：\href{mailto:xxx@qq.com}{250xxxxxxx@qq.com}
        \hspace{3.2em}
        \faPhone \quad 电话：198xxxxxxxx
        \hspace{3.2em}
        \faGithub \quad 主页：\href{https://github.com/Dithob}{github.com/Dithob}
    }
}

\begin{document}

    \begin{tikzpicture}[remember picture, overlay]
        \node[anchor=north, inner sep=0pt](headercontacts) at (current page.north){
            \includegraphics[width=\paperwidth]{images/foot.png}
        };
        \node[anchor=center] at(headercontacts.center){\contact};
    \end{tikzpicture}
    \vspace{-1.15em}

    %%%%%%%%%%%%%%%%%%%%
    % 个人信息与教育背景
    %%%%%%%%%%%%%%%%%%%%

    \begin{minipage}[t]{0.79\textwidth}
        \section[个人信息]{\makebox[\iconwidth][c]{\color{primarycolor}{\faAddressCard}}\quad 个人信息}
        \begin{tabular*}{\textwidth}{@{\extracolsep{\fill}}lll}
            \textbf{姓\qquad 名}：无绝 & \textbf{出生年月}：2002.10 & \textbf{籍\qquad 贯}：湖北咸宁 \\
            \textbf{学\qquad 校}：重庆邮电大学 & \textbf{英语水平}：CET-6 & \textbf{求职方向}：测试开发
        \end{tabular*}

        \section[教育背景]{\makebox[\iconwidth][c]{\color{primarycolor}{\faGraduationCap}}\quad 教育背景}
        {\textbf{重庆邮电大学}}，计算机技术，\textit{硕士} \hfill 2024.09--2027.06
        \begin{itemize}
            \item \textbf{荣誉奖项}：特等学业奖学金 $\times$ 1、二等学业奖学金 $\times$ 1
        \end{itemize}

        \vspace{0.2em}
        {\textbf{重庆邮电大学}}，计算机科学与技术，\textit{本科} \hfill 2020.09--2024.07
        \begin{itemize}
            \item \textbf{荣誉奖项}：三等学业奖学金；GPA: 3.28/4.0（专业前 25\%）
        \end{itemize}
    \end{minipage}
    \hfill
    \begin{minipage}[t]{0.14\textwidth}
        \vspace{-0.35em}
        \setlength{\fboxsep}{0pt}
        \doublebox{\includegraphics[width=\linewidth]{images/xjpic.jpg}}
    \end{minipage}
    % \vspace{0.45em}

    %%%%%%%%%%%%%%%%%%%%
    % 专业技能
    %%%%%%%%%%%%%%%%%%%%

    \section[专业技能]{\makebox[\iconwidth][c]{\color{primarycolor}{\faWrench}}\quad 专业技能}
    \begin{itemize}
        \item \textbf{测试开发基础}：熟悉测试流程、用例设计、等价类/边界值/场景法、缺陷生命周期与回归测试；掌握接口测试、Web UI 自动化、性能测试和稳定性测试，能够围绕业务链路设计冒烟、回归和异常场景用例。
        \item \textbf{自动化测试工具}：熟悉 Python、Pytest、Requests、Postman、JMeter/Locust、Allure、Playwright/Selenium；能够封装接口客户端、公共断言、Fixture、参数化数据驱动、失败重试、日志采集、截图/Trace 留存和测试报告。
        \item \textbf{工程与平台能力}：掌握 Linux、Git、Docker、MySQL、Redis、FastAPI/SpringBoot、Vue.js 等基础技术；了解 Jenkins/GitHub Actions、Mock 服务、测试数据构造、接口契约校验、SQL 校验和 CI 质量门禁。
        \item \textbf{AI 应用质量保障}：理解 RAG、Prompt、向量检索、LLM API 集成与模型评估；能够构建 Golden Dataset、Bad Case 回流、幻觉/拒答/格式化输出校验和多轮对话回归测试集。
    \end{itemize}

    % \vspace{0.35em}

    %%%%%%%%%%%%%%%%%%%%
    % 实习经历
    %%%%%%%%%%%%%%%%%%%%

    \section[实习经历]{\makebox[\iconwidth][c]{\color{primarycolor}{\faBriefcase}}\quad 实习经历}

    \projectheading{重庆啄木鸟网络科技有限公司}{算法实习生（AI 应用测试与交付）}{2025.06--2025.09}
    \begin{itemize}
        \item \textbf{核心职责}：参与大模型应用研发与交付验证，负责模型服务接口联调、Prompt/RAG 流程测试、功能冒烟与回归用例设计、接口异常排查和线上 Bad Case 分析，推动模型回复质量与服务稳定性优化。
        \item \textbf{小啡 GPT 智能回复模型}：围绕家电客服多轮对话构建意图识别、槽位抽取、知识命中、指令遵循和结构化输出测试集；使用脚本批量调用模型 API，对回复内容进行关键词、JSON Schema、相似度和规则断言校验，沉淀 100+ 条对话回归样例，定位上下文丢失、知识召回错误和格式漂移等问题。
        \item \textbf{智能家电识别与自动入库系统}：针对 OCR/视觉大模型识别、联网补全、向量去重和入库链路设计接口测试与数据质量校验；使用 Postman/Python 构造正常、缺字段、重复商品、低置信度识别等场景，并通过数据库断言验证字段完整性、去重结果和入库状态，提升交付前缺陷发现效率。
    \end{itemize}

    % \vspace{0.35em}

    %%%%%%%%%%%%%%%%%%%%
    % 项目经历
    %%%%%%%%%%%%%%%%%%%%

    \section[项目经历]{\makebox[\iconwidth][c]{\color{primarycolor}{\faProjectDiagram}}\quad 项目经历}

    \projectheading{接口自动化测试框架}{核心开发}{2025.10--2025.12}
    \begin{itemize}
        \item \textbf{项目目标}：参考 Pytest、Schemathesis 与 Allure，面向登录、用户、权限、商品、订单等接口搭建可复用自动化框架；技术栈为 Python、Pytest、Requests、YAML/JSON、Allure、MySQL、Docker、Jenkins/GitHub Actions。
        \item \textbf{框架设计}：按配置层、API 客户端层、业务关键字层、用例层和报告层拆分；封装环境切换、Token 自动刷新、统一请求日志、公共断言、数据库校验、Fixture 和参数化数据驱动，支持冒烟/回归/按模块标签选择执行。
        \item \textbf{契约与异常测试}：基于 OpenAPI/Swagger 文档生成契约用例，自动构造边界值、非法枚举、缺失必填字段和嵌套对象异常数据，覆盖状态码、响应 Schema、幂等性和业务错误码，输出失败接口的请求参数、响应体和可复现 curl。
        \item \textbf{质量闭环}：集成 Allure/JUnit XML 报告与 CI 流水线，代码提交后自动执行冒烟集合，失败时保留日志、接口耗时和失败历史；通过报告趋势观察高频失败接口和慢接口，为缺陷定位与回归范围选择提供依据。
    \end{itemize}

    \vspace{0.25em}

    \projectheading{持续测试平台与性能压测}{主要开发}{2026.01--2026.04}
    \begin{itemize}
        \item \textbf{平台能力}：参考 MeterSphere、Playwright 与 Locust，使用 Vue.js、FastAPI、MySQL、Redis、Docker 搭建轻量测试平台，支持项目管理、用例维护、测试计划、任务记录、缺陷备注和报告归档，统一管理接口、UI 与性能测试结果。
        \item \textbf{UI 自动化}：基于 Playwright 落地 PO 模式，封装表单提交、文件上传等核心业务用例；通过内置 智能等待 与 语义化定位器 降低脚本脆弱性；集成 Trace 、失败截图与视频录制，支持全链路回放快速定位 DOM 变更、网络异常与 终端报错。
        \item \textbf{性能测试}：使用 JMeter/Locust 编写登录鉴权、商品查询、订单提交等压测场景，设计阶梯加压、稳定性运行和峰值压力测试；统计 QPS、平均响应时间、P95、错误率和慢接口列表，并结合日志定位数据库查询、缓存命中和接口超时问题。
        \item \textbf{工程落地}：将 Pytest、Playwright 与 Locust 执行入口封装为平台任务，通过 Redis 队列异步调度并回收结果；报告页聚合通过率、失败原因、耗时分布和历史趋势，形成“用例设计-自动执行-报告分析-缺陷回归”的闭环。
    \end{itemize}

    % \vspace{0.35em}

    %%%%%%%%%%%%%%%%%%%%
    % 奖项证书
    %%%%%%%%%%%%%%%%%%%%

    \section[在校成果]{\makebox[\iconwidth][c]{\color{primarycolor}{\faAward}}\quad 在校成果}
    \begin{itemize}
        \item \textbf{竞赛}：2025 “互联网+”大学生创新创业大赛银奖；2025 中国国际大学生创新大赛产业赛道校级二等奖；2025 首届重庆市 AI 大模型创新应用大赛省级三等奖
        \item \textbf{论文}：《LHAF: A Lightweight Hierarchical Adaptive Framework for LLM-based Multimodal Emotion Recognition in Conversations》（SCI I 区在投）
    \end{itemize}

\end{document}
